\section{Conclusion}
\label{sec:conclusion}

Sharding makes a single committee cheap to corrupt and bribery makes corruption purchasable, but
neither fact is consequential while accountable BFT can slash and revert in time. We have argued that
the real question is a race, and we have given it a model: a bribe-cost expression in which the price
of equivocation is set by the probability $q$ that the fraud proof is enforced before the validator's
stake becomes unslashable, coupled to an accountability-race condition that factors viability into a
window gate $\Tsettle<\Tacc$ and an economic gate $\Eval>B+\Cop$. We were precise that the $F{+}1$
equivocators are necessary but not sufficient---success additionally requires an honest-vote split
captured by $\psucc$---and that ``sub-threshold'' is a network-wide, not a committee-local, statement.
The accountability latency $\Tacc$ emerges as a master variable that both widens the window and
cheapens the bribe---``cheap to corrupt'' and ``easy to cash out'' are two symptoms of one
quantity---and we proved that finality-gating cross-shard acceptance closes the attack
unconditionally. Extending the model to epoch boundaries, we showed that reconfiguration does not
remove the race but relocates it: a safe-handoff theorem gives the four conditions under which old-epoch
artifacts are immune, and a residual-vulnerability theorem shows that, absent them, last-of-epoch
attacks strictly dominate. The framework is constructive in both directions---it tells an attacker
exactly when sub-threshold committee bribery pays, and a designer exactly what to enforce so that it
never does---and it is reproducible: a seeded simulator and a local proof-of-concept regenerate every
figure and table, with the harness reproducing the closed-form success probability to within
Monte-Carlo error. The decisive next step is to measure the two load-bearing quantities, $q$ and
$\Tacc$, on one concrete deployed sharded system, and thereby exhibit---or rule out---a real design
with an open window and profitable economics.
